% TOP_PROBLEM: {"problemId": "problem.the-fontaine-mazur-conjecture-on-geometric-galois-representations", "canonicalTitle": "The Fontaine-Mazur Conjecture on Geometric Galois Representations", "releaseRank": 47, "primaryDomain": "number_theory_arithmetic_geometry", "primaryDomainLabel": "Number theory & arithmetic geometry", "displayStatus": "open_with_solved_subcases", "releaseStatus": "open_with_solved_subcases"}
% !TeX program = lualatex
% Definition and sources only; this file is not an LLM solution attempt.
\documentclass[11pt]{article}
\usepackage[margin=25mm]{geometry}
\usepackage{fontspec}
\setmainfont{Latin Modern Roman}
\usepackage{amsmath,amssymb,mathtools}
\usepackage{unicode-math}
\setmathfont{Latin Modern Math}
\usepackage{xurl,hyperref}
\hypersetup{colorlinks=true,urlcolor=blue}
\setcounter{secnumdepth}{0}
\setlength{\parindent}{0pt}
\setlength{\parskip}{5pt}
\setlength{\emergencystretch}{3em}
\begin{document}
\section*{\texorpdfstring{The Fontaine-Mazur Conjecture on Geometric Galois Representations}{The Fontaine-Mazur Conjecture on Geometric Galois Representations}}\label{47}
\subsection{Definitions and mathematical statement}
Let $G_K=\operatorname{Gal}(\bar K/K)$ and $\rho:G_K\to\operatorname{GL}(V)$ be an irreducible continuous representation over a finite extension of ${\mathbb{Q}}_\ell$. Being unramified outside a finite set means inertia acts trivially at all remaining finite places. At $v\mid\ell$, being de Rham means the de Rham period module has the full dimension of $V$. The assertion is that
\[
 V\text{ is a subquotient of }H^i_{\mathrm{\acute et}}(X_{\bar K},{\mathbb{Q}}_\ell)\otimes{\mathbb{Q}}_\ell(m)
\]
with coefficients extended to those of $V$, for some smooth projective $X/K$, $i\ge0$, and $m\in{\mathbb{Z}}$. Arbitrary Tate twists are allowed; no nonnegative Hodge--Tate-weight condition is imposed.
\par\smallskip\textit{Formulation and concept references:} \href{https://math.berkeley.edu/~swshin/FieldRat.pdf}{[S1]}.

\subsection{Short English statement}
Must every irreducible Galois representation with the specified arithmetic and local geometric properties arise from algebraic geometry?
\subsection{Sources}
\begin{itemize}
\item[S1]Sug Woo Shin and Nicolas Templier, On fields of rationality for automorphic representations, author version dated 2 July 2014.
\url{https://math.berkeley.edu/~swshin/FieldRat.pdf}.
\textit{Formulation source.}
\item[S2]James Newton, Modularity of Galois Representations and Langlands Functoriality.
\url{https://link.springer.com/article/10.1007/s41745-022-00305-0}.
\textit{Further reference; not a proof endorsement.}
\item[S3]Jack A. Thorne, Towards the Fontaine–Mazur conjecture for GL(2).
\url{https://arxiv.org/html/2608.07186v1}.
\textit{Further reference; not a proof endorsement.}
\end{itemize}
\end{document}
